\documentclass[11pt]{article}

\usepackage[margin=1in]{geometry}
\usepackage{amsmath, amssymb}
\usepackage{xcolor}
\usepackage[skins,breakable]{tcolorbox}
\usepackage{hyperref}

\hypersetup{
    colorlinks=true,
    linkcolor=blue!60!black,
    urlcolor=blue!60!black,
    citecolor=blue!60!black
}

\newcommand{\SymPy}{SymPy}

% Public artifact repository; \bugbundle links a bug's bundle folder into it.
% TODO: set the public artifact repository URL before release.
\newcommand{\artifactrepo}{https://github.com/TODO-org/TODO-repo}
\newcommand{\bugbundle}[1]{\href{\artifactrepo/tree/main/bug-report/#1}{\nolinkurl{bug-report/#1/}}}

\newtcolorbox{counterexamplebox}{
  enhanced,
  colback=yellow!5,
  colframe=orange!70!black,
  arc=2mm,
  boxrule=0.8pt,
  left=4mm,
  right=4mm,
  top=3mm,
  bottom=3mm,
  before=\par\medskip\noindent,
  after=\par\medskip,
  before upper=\raggedright
}

\title{SymPy Bug Card}
\author{}
\date{}

\begin{document}

\maketitle

\section*{Bug 33: solveset Represents an Unsatisfiable Tangent Equation by Infinities}

\begin{counterexamplebox}
\textbf{Task.} Solve the complex trigonometric equation
\[
\tan(x)=i,\qquad x\in\mathbb{C}.
\]

\medskip
\textbf{SymPy's answer.}
\[
\operatorname{solveset}(\tan(x)=i,\ x,\ \mathbb{C})
= \operatorname{ImageSet}(\Lambda(\_n,\ \_n\pi+oo\,i),\mathbb{Z}).
\]

\medskip
\textbf{Correct answer.}
\[
\varnothing.
\]

\medskip
\textbf{Explanation.} For finite \(z\),
\[
\tan(z)=-i\,\frac{\exp(2iz)-1}{\exp(2iz)+1}.
\]
Setting \(\tan(z)=i\) forces \(\exp(2iz)=0\), impossible for a finite complex
argument. The returned \(oo\,i\)-valued image set is therefore not a finite
complex solution family.

\medskip
\textbf{Likely root cause.} In \nolinkurl{sympy/solvers/solveset.py},
\texttt{\_invert\_trig\_hyp\_complex} constructs the tangent solution pattern
\(n\pi+\arctan(a)\) without rejecting non-finite inverse values; here
\(\arctan(i)\) evaluates to \(oo\,i\). This is a new instance of a known failure
mode.

\medskip
\textbf{Artifact (GitHub).} \bugbundle{bug-033-solveset-represents-an-unsatisfiable-tangent-equation-by-inf}
\end{counterexamplebox}

\end{document}
